\chapter{Active Circuits, Feedback, and Oscillators}
\label{ch:active}
\section{Why Active Circuits Appear}
The Extra exam includes amplifiers, op-amps, feedback, oscillators, mixers, and
regulated power supplies. A full electronics text is not needed here; the goal is
to connect the circuit formulas to system behavior---and, in this chapter, to the
classical feedback ideas that a controls background already contains. The device
details (bias points, transistor models, tube parameters) live in the ARRL
manual; what follows is the systems view that ties them together. The feedback
groundwork of \cref{ch:feedback}---closed-loop gain \(A/(1+A\beta)\), the
large-loop-gain limit \(A_{\text{cl}}\to1/\beta\), pole motion under feedback, and
the gain and phase margins---is the backbone of this chapter. Rather than restate
those results, this chapter puts them to work on real devices, so it is worth having
\cref{ch:feedback} fresh in mind.

\section{Amplifier Gain and Classes}
Voltage gain is \(A_v=V_o/V_i\); in decibels \(A_{v,\dB}=20\log_{10}|A_v|\). Power
gain is \(G_{\dB}=10\log_{10}(P_o/P_i)\).

Amplifier \emph{class} describes how much of the input cycle the active device
conducts, which trades linearity against efficiency:
\begin{description}
\item[Class A] conducts the full \(360^\circ\); most linear, least efficient
(theoretical max \SI{50}{\percent} transformer-coupled, near \SI{25}{\percent}
resistively coupled).
\item[Class B / AB] Class B conducts half the cycle (\(180^\circ\)); Class AB
somewhat more than half. Higher efficiency than Class A, with some crossover
distortion; common in linear single-sideband (SSB) power amplifiers.
\item[Class C] conducts less than half the cycle; efficient but nonlinear, used
where a tuned output tank restores the sinusoid (as in frequency-modulation, FM,
and continuous-wave Morse, CW, transmitters).
\end{description}
Class C is nonlinear on purpose: the output tank is a high-\(Q\) resonator (a
pole pair near the \(\jj\omega\) axis) that rings at the fundamental and rejects
the harmonics the clipped current injects. The same nonlinearity that spoils
linearity is what makes mixers and frequency multipliers work.

\section{Op-Amp Ideals}
An ideal op-amp has very high input impedance, very low output impedance, and
very large open-loop gain \(A(s)\).

\begin{physicalbox}
\textbf{How the op-amp is modeled.} An op-amp is \emph{not} a reactive element
like an inductor or capacitor---it stores no energy and contributes no pole of its
own in the ideal picture. It is modeled as a \textbf{dependent source}:
specifically a voltage-controlled voltage source (VCVS) whose output is
\(A(s)\,(v_+-v_-)\), the open-loop gain times the difference of the two input
voltages. The idealizations attach to that source: infinite input impedance means
no current flows into the \(\pm\) terminals; zero output impedance means the
controlled source drives the output stiffly; and \(A(s)\) is enormous (and rolls
off with frequency, \cref{sec:gbw}). Every op-amp result below is bookkeeping on
this one dependent-source model plus the feedback wired around it.
\end{physicalbox}

Negative feedback makes practical closed-loop gains predictable:
\[
\text{inverting:}\quad \frac{V_o}{V_i}=-\frac{R_f}{R_{in}},
\qquad
\text{non-inverting:}\quad \frac{V_o}{V_i}=1+\frac{R_f}{R_g}.
\]
These come from the two ideal rules---no current into the inputs, and the inputs
driven to equal voltage by feedback---but both rules are consequences of a deeper
statement: the loop gain is large. The inverting configuration
(\cref{fig:opamp-inv}) makes this concrete: with the non-inverting input grounded
and the inverting input held at a virtual ground by feedback, the input current
\(V_i/R_{in}\) must flow on through \(R_f\), giving \(V_o=-R_f\,(V_i/R_{in})\).

\begin{figure}[htbp]
\centering
\begin{circuitikz}[scale=1]
\node[op amp,yscale=-1] (oa) at (0,0) {};
\draw (oa.-) to[short] ++(-0.7,0) coordinate(n);
\draw (n) to[R,l_=\(R_{in}\)] ++(-2.2,0) node[left]{\(V_i\)};
\draw (oa.+) to[short] ++(-0.7,0) node[ground]{};
\draw (oa.out) to[short,-o] ++(0.9,0) coordinate(vo) node[right]{\(V_o\)};
\draw (n) -- ++(0,1.6) coordinate(tl) to[R,l=\(R_f\)] (tl-|vo) -- (vo);
\end{circuitikz}
\caption{Inverting op-amp amplifier. Feedback holds the inverting input at a
virtual ground, so \(V_o/V_i=-R_f/R_{in}\)---set by the resistor ratio, not the
op-amp.}
\label{fig:opamp-inv}
\end{figure}

\section{Feedback as a Loop}
Model the amplifier as a forward path \(A(s)\) with a fraction \(\beta\) of the
output fed back and subtracted at the input.

\begin{figure}[htbp]
\centering
\begin{tikzpicture}[>=Latex,node distance=1.1cm,scale=1]
\node (in) at (0,0) {\(V_i\)};
\node[draw,circle,inner sep=1.4pt] (sum) at (1.3,0) {\(\Sigma\)};
\node[draw,minimum width=1.3cm,minimum height=0.8cm] (A) at (3.6,0) {\(A(s)\)};
\node (out) at (6.3,0) {\(V_o\)};
\node[draw,minimum width=1.1cm,minimum height=0.7cm] (B) at (3.6,-1.5) {\(\beta\)};
\draw[->] (in) -- (sum);
\draw[->] (sum) -- (A);
\draw[->] (A) -- (out);
\draw[->] (5.3,0) |- (B);
\draw[->] (B) -| (sum);
\node at (0.95,0.3) {\small \(+\)};
\node at (1.62,-0.28) {\small \(-\)};
\end{tikzpicture}
\caption{Standard negative-feedback loop: the reference \(V_i\) enters the summing
junction with \(+\) and the feedback \(\beta V_o\) is subtracted (\(-\)). The loop
gain is \(L(s)=A(s)\beta\).}
\label{fig:feedback-loop}
\end{figure}

The closed-loop gain is
\[
A_{\text{cl}}(s)=\frac{A(s)}{1+A(s)\beta}
=\frac{A(s)}{1+L(s)},
\qquad
L(s)=A(s)\beta .
\]
When the loop gain \(|L|\gg1\), \(A_{\text{cl}}\to 1/\beta\): the gain is set by
the passive feedback network, not by the messy device parameters. That single
fact is why the inverting and non-inverting formulas above depend only on
resistor ratios.

\begin{controlsbox}
\textbf{What this buys.} By the sensitivity result of \cref{ch:feedback},
\(\dd A_{\text{cl}}/A_{\text{cl}}=(\dd A/A)/(1+L)\): the loop divides the device's
imperfections by \(1+L\). For an op-amp that is the whole design philosophy---spend
the enormous surplus gain \(A\) to buy immunity from device variation, temperature,
and distortion, and let two resistors set the answer.
\end{controlsbox}

\section{Gain--Bandwidth Product}
\label{sec:gbw}
An internally compensated op-amp is deliberately built so that one pole dominates,
\[
A(s)=\frac{A_0}{1+s/\omega_a},
\]
with an enormous DC gain \(A_0\) (often \(10^5\)) and a very low open-loop corner
\(\omega_a\) (often a few hertz). That is precisely the first-order plant analyzed
in \cref{ch:feedback}, so its result transfers unchanged: closing the loop with
constant \(\beta\) divides the DC gain by \((1+A_0\beta)\) and multiplies the corner
by the same factor, leaving the product invariant,
\[
(\text{gain})\times(\text{bandwidth})=A_0\,\omega_a=\text{constant}.
\]
For amplifiers this invariant has a name and a datasheet value: the
\textbf{gain--bandwidth product}, quoted as the \emph{unity-gain frequency}
\(f_T=A_0f_a\)---the frequency at which the op-amp can no longer provide any gain at
all. Every closed-loop configuration you build with that part lands on the same
diagonal (\cref{fig:opamp-gbw}), so choosing a closed-loop gain \emph{is} choosing a
bandwidth: doubling one halves the other.

\begin{figure}[htbp]
\centering
\includegraphics[width=0.86\textwidth]{opamp_gbw_bode}
\caption{Gain--bandwidth product of a single-pole op-amp. Each closed-loop gain
(set by \(1/\beta\)) is flat until it meets the open-loop skirt, then rolls off
along it, so every closed-loop corner lands on the same unity-gain frequency
\(f_T=A_0 f_a\). Trading gain for bandwidth slides along one fixed diagonal.}
\label{fig:opamp-gbw}
\end{figure}

\begin{exambox}
Gain--bandwidth questions are the practical face of this: an op-amp with
\(f_T=\SI{1}{\mega\hertz}\) configured for a gain of \(100\) has roughly
\SI{10}{\kilo\hertz} of bandwidth, not \SI{1}{\mega\hertz}. If a stage needs both
high gain and wide bandwidth, it needs either a faster part or two cascaded stages
of modest gain.
\end{exambox}

As \cref{ch:feedback} showed, this single-pole model can never be destabilized by
more feedback---the pole only slides left. Instability requires the extra phase lag
of additional poles, which every real amplifier has.

\section{Stability, Phase Margin, and Gain Margin}
Real amplifiers have more than one pole. With two or more poles the loop gain
\(L(\jj\omega)\) accumulates phase lag, and if the phase reaches \(-180^\circ\)
while \(|L|\ge 1\), the subtracting junction effectively adds---the feedback has
become positive at that frequency and the closed-loop poles move toward, or
into, the right half-plane.

The gain and phase margins of \cref{ch:feedback} are exactly how amplifier
stability is specified and measured: positive margins keep the closed-loop poles in
the left half-plane, small margins put them near the \(\jj\omega\) axis (ringing,
overshoot, a peaked response), and zero margin puts them on it. An op-amp datasheet's
phase-margin figure is that number for the part's own internal loop.

The same story reads on the s-plane as a root locus (\cref{fig:opamp-rootlocus}):
as loop gain rises, the closed-loop poles of a multi-pole amplifier migrate, and
with three poles two of them curve toward and cross the \(\jj\omega\) axis---the
moment small margins become ringing and then sustained oscillation.

\begin{figure}[htbp]
\centering
\includegraphics[width=0.6\textwidth]{opamp_rootlocus}
\caption{Root locus of a three-pole amplifier as loop gain \(K\) increases. Two
closed-loop poles leave the real axis and cross into the right half-plane at
\(K=60\): beyond that the amplifier oscillates. A single-pole amplifier (one
branch on the real axis) can never do this---instability needs the extra phase
lag of additional poles.}
\label{fig:opamp-rootlocus}
\end{figure}

Compensation (deliberately lowering the first pole, as in an internally
compensated op-amp) restores phase margin by making the loop gain fall through
\SI{0}{\decibel} before the second pole adds its lag. It buys stability with
bandwidth---the same trade, again.

\section{Oscillators: Poles Placed on the Axis}
An oscillator is a feedback loop designed to sit exactly on the stability
boundary. The closed-loop poles are the roots of \(1+L(s)=0\), so they reach the
\(\jj\omega\) axis when
\[
L(\jj\omega_{\text{osc}})=-1,
\qquad\text{that is}\qquad
|L(\jj\omega_{\text{osc}})|=1,\quad
\angle L(\jj\omega_{\text{osc}})=-180^\circ,
\]
the Nyquist \(-1\) point. This is the \emph{Barkhausen criterion}. Its familiar
``unit magnitude, zero net phase'' wording refers to the round-trip transmission
\(T=-L\), which already includes the \(-1\) of the subtracting junction; setting
\(T=1\) is the same statement as \(L=-1\). In s-plane language the closed-loop
poles are placed on the \(\jj\omega\) axis at \(\pm\jj\omega_{\text{osc}}\)---the
marginally stable case from the s-plane chapter. At start-up the design makes \(|L|>1\) so the pole pair sits slightly in
the right half-plane and the amplitude grows; a nonlinearity (device saturation,
a limiter, a lamp, or automatic gain control, AGC) then reduces the effective gain until \(|L|=1\) and the
poles settle onto the axis at a stable amplitude.

\Cref{fig:oscillator-loopgain} shows the condition on a loop-gain Bode plot: the
design places the frequency where the phase reaches \(-180^\circ\) exactly at the
\(\SI{0}{\decibel}\) (\(|L|=1\)) crossing, so \(L=-1\) is met at one frequency.

\begin{figure}[htbp]
\centering
\includegraphics[width=0.82\textwidth]{oscillator_loopgain}
\caption{Barkhausen condition for a three-pole loop
\(L(s)=K/(1+s/\omega_0)^3\) with \(K=8\). The magnitude passes through
\SI{0}{\decibel} at the same frequency \(\omega_{\mathrm{osc}}=\sqrt3\,\omega_0\)
where the phase reaches \(-180^\circ\): loop gain unity with \(180^\circ\) of lag
is \(L=-1\), so the closed-loop poles sit on the \(\jj\omega\) axis.}
\label{fig:oscillator-loopgain}
\end{figure}

The frequency-setting network sets where those poles land:
\begin{description}
\item[LC oscillators] (Colpitts, Hartley) use an LC tank; \(\omega_{\text{osc}}
\approx 1/\sqrt{LC}\). Higher tank \(Q\) means poles closer to the ideal axis
location and lower phase noise.
\item[Crystal oscillators] replace the tank with a quartz resonator of very high
\(Q\), giving excellent frequency stability.
\item[RC phase-shift / Wien-bridge] use RC networks to reach the required loop
phase; common at audio frequencies.
\end{description}

\begin{physicalbox}
Phase noise is the frequency-domain fingerprint of how firmly the poles are
pinned to the axis. A high-\(Q\) resonator holds the pole pair tightly at
\(\pm\jj\omega_{\text{osc}}\), so the spectrum is a sharp line; a low-\(Q\) loop
lets the poles wander, smearing the line. This is why crystal and high-\(Q\) LC
oscillators are quieter.
\end{physicalbox}

\section{Active Filters and Their Limits}
Op-amps can implement active RC filters that synthesize the second-order pole
pairs of \cref{ch:filters} without inductors, by using feedback to create complex
poles from resistors and capacitors. The Sallen--Key low-pass
(\cref{fig:sallenkey}) is the classic example: two RC stages around a unity-gain
buffer place a complex pole pair whose \(\wzero=1/\sqrt{R_1R_2C_1C_2}\) and whose
\(Q\) is set by the component ratios---no inductor required.

\begin{figure}[htbp]
\centering
\begin{circuitikz}[scale=1]
\node[op amp] (oa) at (7,0) {};
\draw (0,0) node[left]{\(V_{\mathrm{in}}\)} to[R,l=\(R_1\)] (2,0) coordinate(a)
      to[R,l=\(R_2\)] (4,0) coordinate(b);
\draw (b) -| (oa.+);
\draw (b) to[C,l_=\(C_1\)] ($(b)+(0,-1.8)$) node[ground]{};
\draw (oa.out) to[short,-o] (9,0) coordinate(vo) node[right]{\(V_o\)};
% unity-gain feedback from inverting input up and over to output
\draw (oa.-) -- ++(0,1.2) -| (vo);
% C2 from node a up and over to the output
\draw (a) -- ++(0,2.5) coordinate(ct) to[C,l=\(C_2\)] (ct-|vo) -- (vo);
\end{circuitikz}
\caption{Sallen--Key second-order low-pass. Two RC stages around a unity-gain
op-amp buffer create a complex pole pair without an inductor;
\(\wzero=1/\sqrt{R_1R_2C_1C_2}\).}
\label{fig:sallenkey}
\end{figure}

The catch is that the same finite gain--bandwidth product limits them: near the
amplifier's own gain-crossover the ideal-op-amp assumption fails, so an active
filter that behaves at audio may be useless at RF, where LC or crystal filters
take over.

\section{Power Supplies}
Rectifiers, filter capacitors, bleeder resistors, regulators, and switching
converters reuse the elements from earlier chapters. A linear regulator is itself
a feedback loop---it compares the output to a reference and drives the pass
element to null the error, with the same stability considerations as any feedback
system. Equal-value resistors across series filter capacitors equalize the shared
voltage and discharge stored energy after power is removed.

\begin{exambox}
For active-circuit questions, keep three layers separate: small-signal gain,
power handling, and stability. Feedback questions almost always reduce to one
idea---\emph{loop gain}: large loop gain buys predictable gain, wide bandwidth,
and low distortion; too much phase lag at unity loop gain costs stability; and an
oscillator is the deliberate limiting case with the poles parked on the
\(\jj\omega\) axis.
\end{exambox}
